%2multibyte Version: 5.50.0.2960 CodePage: 65001
%\newtheorem{definition}[definicion]{Definici�n}
%\newtheorem{example}[ejemplo]{Ejemplo}


\documentclass[handout,9pt]{beamer}
%%%%%%%%%%%%%%%%%%%%%%%%%%%%%%%%%%%%%%%%%%%%%%%%%%%%%%%%%%%%%%%%%%%%%%%%%%%%%%%%%%%%%%%%%%%%%%%%%%%%%%%%%%%%%%%%%%%%%%%%%%%%%%%%%%%%%%%%%%%%%%%%%%%%%%%%%%%%%%%%%%%%%%%%%%%%%%%%%%%%%%%%%%%%%%%%%%%%%%%%%%%%%%%%%%%%%%%%%%%%%%%%%%%%%%%%%%%%%%%%%%%%%%%%%%%%
\usepackage{amsmath}
\usepackage{mathpazo}
\usepackage{hyperref}
\usepackage{multimedia}
\usepackage{times}
\usepackage[T1]{fontenc}

\setcounter{MaxMatrixCols}{10}
%TCIDATA{OutputFilter=LATEX.DLL}
%TCIDATA{Version=5.50.0.2960}
%TCIDATA{Codepage=65001}
%TCIDATA{<META NAME="SaveForMode" CONTENT="1">}
%TCIDATA{BibliographyScheme=Manual}
%TCIDATA{LastRevised=Tuesday, May 16, 2023 21:34:10}
%TCIDATA{<META NAME="GraphicsSave" CONTENT="32">}

\newtheorem{definicion}{Definici\'{o}n}
\newtheorem{ejemplo}{Ejemplo}
\newtheorem{ejemplos}{Ejemplos}
\newenvironment{stepenumerate}{\begin{enumerate}[<+->]}{\end{enumerate}}
\newenvironment{stepitemize}{\begin{itemize}[<+->]}{\end{itemize} }
\newenvironment{stepenumeratewithalert}{\begin{enumerate}[<+-| alert@+>]}{\end{enumerate}}
\newenvironment{stepitemizewithalert}{\begin{itemize}[<+-| alert@+>]}{\end{itemize} }
\usetheme{Rochester}
\usecolortheme[RGB={64,63,152}]{structure}
\usefonttheme{professionalfonts}
\input{tcilatex}
\begin{document}

\title[MAT-015]{Principios Matem\'{a}ticos}
\subtitle{Unidad III. Ecuaciones de Primer y Segundo Grado y sus Aplicaciones\\
Clase: Aplicaciones de las Ecuaciones de primer grado}
\author[lgonzalez@santotomas.cl]{Profesor Luis Tulio Gonz\'{a}lez Alcaino \\
%EndAName
Magister en Matem\'{a}tica}
\institute[UST]{Universidad Santo Tom\'{a}s\\
Departamento Ciencias B\'{a}sicas - Talca }
\date{}
\maketitle

%TCIMACRO{\TeXButton{BeginFrame}{\begin{frame}}}%
%BeginExpansion
\begin{frame}%
%EndExpansion

\QTR{frametitle}{Contenidos de la clase}

\begin{itemize}
\item Aplicaciones ecuaciones de primer grado

\item Ejercicios
\end{itemize}

%TCIMACRO{\TeXButton{EndFrame}{\end{frame}}}%
%BeginExpansion
\end{frame}%
%EndExpansion

%TCIMACRO{\TeXButton{BeginFrame}{\begin{frame}}}%
%BeginExpansion
\begin{frame}%
%EndExpansion

\QTR{frametitle}{Problemas con ecuaciones de primer grado}

En la mayor\'{\i}a de los casos, para resolver problemas pr\'{a}cticos deben

traducirse las relaciones establecidas a s\'{\i}mbolos matem\'{a}ticos. Esto
es

conocido como modelaci\'{o}n. No existe un m\'{e}todo o regla general para

resolver problemas de aplicaci\'{o}n, sin embargo, existen algunos pasos a

seguir que con frecuencia son \'{u}tiles para obtener la soluci\'{o}n:

\begin{itemize}
\item Leer detenidamente el problema.

\item Hacer un esquema o dibujo, si tiene sentido hacerlo.

\item Representar el(los) t\'{e}rmino(s) desconocido(s) por medio de una(s)
variable(s).

\item Si es posible, representar todas las dem\'{a}s cantidades en t\'{e}%
rminos de la o las variables utilizadas.

\item Expresar la situaci\'{o}n descrita en el problema en t\'{e}rminos matem%
\'{a}ticos,

\item es decir, plantear la(s) ecuaci\'{o}n(es).

\item Resolver la(s) ecuaci\'{o}n(es) planteadas.

\item Verificar que la soluciones concuerde con las condiciones planteadas
en el problema, analizar.

\item Dar respuesta al problema planteado.
\end{itemize}

%TCIMACRO{\TeXButton{EndFrame}{\end{frame}}}%
%BeginExpansion
\end{frame}%
%EndExpansion

%TCIMACRO{\TeXButton{BeginFrame}{\begin{frame}}}%
%BeginExpansion
\begin{frame}%
%EndExpansion

\QTR{frametitle}{Ejemplos de problemas con ecuaciones de primer grado}

%TCIMACRO{\TeXButton{BeginBlok}{\begin{block}{Ejemplo de aplicaci\'{o}n}}}%
%BeginExpansion
\begin{block}{Ejemplo de aplicaci\'{o}n}%
%EndExpansion

\textbf{Ejemplo 1.} Una empresa agr\'{\i}cola export\'{o} s\'{o}lo un cuarto
de su producci\'{o}n

de duraznos y 2/3 de la producci\'{o}n se consumi\'{o} en la Regi\'{o}n del
Maule. Si

este consumo fue de 240 toneladas, determine la producci\'{o}n de duraznos

que se export\'{o}.

%TCIMACRO{\TeXButton{EndBlok}{\end{block}}}%
%BeginExpansion
\end{block}%
%EndExpansion

%TCIMACRO{\TeXButton{EndFrame}{\end{frame}}}%
%BeginExpansion
\end{frame}%
%EndExpansion

%TCIMACRO{\TeXButton{BeginFrame}{\begin{frame}}}%
%BeginExpansion
\begin{frame}%
%EndExpansion

\QTR{frametitle}{Ejemplos de problemas con ecuaciones de primer grado}

%TCIMACRO{\TeXButton{BeginBlok}{\begin{block}{Ejemplo de aplicaci\'{o}n}}}%
%BeginExpansion
\begin{block}{Ejemplo de aplicaci\'{o}n}%
%EndExpansion

\textbf{Ejemplo 2.} En una cl\'{\i}nica hay 5 adultos mas que ni\~{n}os y 3
ni\~{n}os m\'{a}s que ancianos. Si en total hay 50 pacientes, 
\textquestiondown Cu\'{a}ntos adultos, ni\~{n}os y ancianos hay?

%TCIMACRO{\TeXButton{EndBlok}{\end{block}}}%
%BeginExpansion
\end{block}%
%EndExpansion

%TCIMACRO{\TeXButton{EndFrame}{\end{frame}}}%
%BeginExpansion
\end{frame}%
%EndExpansion

%TCIMACRO{\TeXButton{BeginFrame}{\begin{frame}}}%
%BeginExpansion
\begin{frame}%
%EndExpansion

\QTR{frametitle}{Ejemplos de problemas con ecuaciones de primer grado}

%TCIMACRO{\TeXButton{BeginBlok}{\begin{block}{Ejemplo de aplicaci\'{o}n}}}%
%BeginExpansion
\begin{block}{Ejemplo de aplicaci\'{o}n}%
%EndExpansion

\textbf{Ejemplo 3.} La edad de Pedro es el doble de la edad de Mar\'{\i}a.
Si en cinco

a\~{n}os mas la suma de sus edades ser\'{a} de 43 a\~{n}os. 
\textquestiondown Qu\'{e} edad tienen

actualmente Pedro y Mar\'{\i}a?

%TCIMACRO{\TeXButton{EndBlok}{\end{block}}}%
%BeginExpansion
\end{block}%
%EndExpansion

%TCIMACRO{\TeXButton{EndFrame}{\end{frame}}}%
%BeginExpansion
\end{frame}%
%EndExpansion

%TCIMACRO{\TeXButton{BeginFrame}{\begin{frame}}}%
%BeginExpansion
\begin{frame}%
%EndExpansion

\QTR{frametitle}{Ejercicios}

\begin{enumerate}
\item En tres d\'{\i}as un estudiante con un trabajo Part Time, recibi\'{o}
un total de \$219.600. Si cada d\'{\i}a recibi\'{o} 4/5 del d\'{\i}a
anterior. \textquestiondown Cu\'{a}nto dinero recibi\'{o} por d\'{\i}a?

\item En un consultorio, un pediatra y un oftalm\'{o}logo atendieron en un d%
\'{\i}a a un total de 42 pacientes entre los dos. Si el pediatra hubiese
atendido a seis personas m\'{a}s y el oftalm\'{o}logo a ocho personas menos,
el pediatra habr\'{\i}a atendido el triple de pacientes que el oftalm\'{o}%
logo. \textquestiondown Cu\'{a}ntos pacientes atendi\'{o} el pediatra?

\item El n\'{u}mero de profesionales del \'{a}rea de la salud que ingresan a
trabajar al sector p\'{u}blico est\'{a} compuesto por 88 personas designados
en dos grupos, A y B; los del grupo A son 20 personas m\'{a}s que los del
grupo B. Al respecto, \textquestiondown cu\'{a}l es el n\'{u}mero de
profesionales del \'{a}rea de la salud del grupo A?
\end{enumerate}

%TCIMACRO{\TeXButton{EndFrame}{\end{frame}}}%
%BeginExpansion
\end{frame}%
%EndExpansion

\end{document}
